% Generated draft: requires mathematical and source review.
\documentclass[11pt]{article}
\usepackage[margin=25mm]{geometry}
\usepackage{fontspec}
\setmainfont{Latin Modern Roman}
\usepackage{amsmath,amssymb,mathtools,unicode-math}
\setmathfont{Latin Modern Math}
\usepackage{enumitem,needspace,xurl,hyperref}
\setlength{\parindent}{0pt}
\setlength{\parskip}{5pt}
\setlength{\emergencystretch}{3em}
\newcommand{\N}{\mathbb{N}}
\newcommand{\Z}{\mathbb{Z}}
\newcommand{\Q}{\mathbb{Q}}
\newcommand{\R}{\mathbb{R}}
\newcommand{\C}{\mathbb{C}}
\newcommand{\F}{\mathbb{F}}
\newcommand{\A}{\mathbb{A}}
\newcommand{\PP}{\mathbb{P}}
\newcommand{\E}{\mathbb{E}}
\newcommand{\eps}{\varepsilon}
\newcommand{\dd}{\,\mathrm{d}}
\begin{document}
\section*{UnsolvedMath ID: 801092}
\textbf{Source problem code:} catalogue entry
\subsection{Definitions and mathematical statement}
A polynomial automorphism of \(\C^d\) is a polynomial map \(f:\C^d\to\C^d\) having a polynomial inverse. The supplied record restricts attention to \emph{regular} polynomial automorphisms, but does not define regularity or the object \(S(f)\). Thus the following transcribes the requested conclusions conditionally on those missing definitions; it is not a complete standalone formulation.

For every regular polynomial automorphism \(f\in\operatorname{Aut}[\C^d]\), with \(S(f)\) understood as the intended complex manifold, establish and investigate the following:
\begin{enumerate}
\item There exist a compact complex manifold \(\bar S(f)\), a holomorphic open embedding \(j:S(f)\hookrightarrow\bar S(f)\), and a divisor \(D\) on \(\bar S(f)\) such that
\[
\bar S(f)\setminus j(S(f))=\operatorname{Supp}(D).
\]
Here an open embedding identifies its domain biholomorphically with an open subset, and \(\operatorname{Supp}(D)\) denotes the underlying analytic hypersurface of the divisor. This expresses the source's set-theoretic assertion that the complement is a divisor; multiplicities are not specified.
\item Any two compact complex manifolds satisfying the preceding compactification requirement for the same \(f\) are bimeromorphic, meaning that there is a meromorphic map between them with a meromorphic inverse. The source does not specify whether this bimeromorphism must restrict to the identity under the two identifications of \(S(f)\).
\item Describe complex objects on \(\bar S(f)\), including analytic subvarieties, vector fields, holomorphic foliations, and positive closed currents, and compute its deformation space.
\end{enumerate}
The dimension range, the construction of \(S(f)\), and the intended meaning of deformation space must be supplied before these requests become precise. The last request also leaves unspecified whether a particular compactification is intended or whether the description should apply to every compactification satisfying the first requirement.
\subsection{Short English statement}
For each regular polynomial automorphism, the problem asks for a compactification of its associated space \(S(f)\) with divisor boundary, uniqueness up to bimeromorphism, and a description of geometric objects and deformations. The supplied record does not define \(S(f)\) or regularity, so this is a conditional statement requiring mathematical review.
\subsection{Sources}
\begin{itemize}
\item Local UnsolvedMath registry, permanent ID \(801092\), source problem code \texttt{catalogue entry}; supplied title: \emph{Henon Maps: Open Problems --- Problem 44}. The registry labels this record \emph{open}; that label is supplied source data, not an independently verified mathematical assessment.
\item Berger and Bedford, \emph{Henon Maps: A List of Open Problems} (2023), Problem 44, as attributed in the supplied record. Supplied source URL: \url{https://arxiv.org/abs/2312.03907}. The background records access on 2026-10-16 and extraction method \texttt{arxiv\_latex}; these are supplied provenance statements, not a claim of independent source consultation.
\end{itemize}
\end{document}
